% generated by GAPDoc2LaTeX from XML source (Frank Luebeck)
\documentclass[a4paper,11pt]{report}

\usepackage[top=37mm,bottom=37mm,left=27mm,right=27mm]{geometry}
\sloppy
\pagestyle{myheadings}
\usepackage{amssymb}
\usepackage[utf8]{inputenc}
\usepackage{makeidx}
\makeindex
\usepackage{color}
\definecolor{FireBrick}{rgb}{0.5812,0.0074,0.0083}
\definecolor{RoyalBlue}{rgb}{0.0236,0.0894,0.6179}
\definecolor{RoyalGreen}{rgb}{0.0236,0.6179,0.0894}
\definecolor{RoyalRed}{rgb}{0.6179,0.0236,0.0894}
\definecolor{LightBlue}{rgb}{0.8544,0.9511,1.0000}
\definecolor{Black}{rgb}{0.0,0.0,0.0}

\definecolor{linkColor}{rgb}{0.0,0.0,0.554}
\definecolor{citeColor}{rgb}{0.0,0.0,0.554}
\definecolor{fileColor}{rgb}{0.0,0.0,0.554}
\definecolor{urlColor}{rgb}{0.0,0.0,0.554}
\definecolor{promptColor}{rgb}{0.0,0.0,0.589}
\definecolor{brkpromptColor}{rgb}{0.589,0.0,0.0}
\definecolor{gapinputColor}{rgb}{0.589,0.0,0.0}
\definecolor{gapoutputColor}{rgb}{0.0,0.0,0.0}

%%  for a long time these were red and blue by default,
%%  now black, but keep variables to overwrite
\definecolor{FuncColor}{rgb}{0.0,0.0,0.0}
%% strange name because of pdflatex bug:
\definecolor{Chapter }{rgb}{0.0,0.0,0.0}
\definecolor{DarkOlive}{rgb}{0.1047,0.2412,0.0064}


\usepackage{fancyvrb}

\usepackage{mathptmx,helvet}
\usepackage[T1]{fontenc}
\usepackage{textcomp}


\usepackage[
            pdftex=true,
            bookmarks=true,        
            a4paper=true,
            pdftitle={Written with GAPDoc},
            pdfcreator={LaTeX with hyperref package / GAPDoc},
            colorlinks=true,
            backref=page,
            breaklinks=true,
            linkcolor=linkColor,
            citecolor=citeColor,
            filecolor=fileColor,
            urlcolor=urlColor,
            pdfpagemode={UseNone}, 
           ]{hyperref}

\newcommand{\maintitlesize}{\fontsize{50}{55}\selectfont}

% write page numbers to a .pnr log file for online help
\newwrite\pagenrlog
\immediate\openout\pagenrlog =\jobname.pnr
\immediate\write\pagenrlog{PAGENRS := [}
\newcommand{\logpage}[1]{\protect\write\pagenrlog{#1, \thepage,}}
%% were never documented, give conflicts with some additional packages

\newcommand{\GAP}{\textsf{GAP}}

%% nicer description environments, allows long labels
\usepackage{enumitem}
\setdescription{style=nextline}

%% depth of toc
\setcounter{tocdepth}{1}





%% command for ColorPrompt style examples
\newcommand{\gapprompt}[1]{\color{promptColor}{\bfseries #1}}
\newcommand{\gapbrkprompt}[1]{\color{brkpromptColor}{\bfseries #1}}
\newcommand{\gapinput}[1]{\color{gapinputColor}{#1}}


\begin{document}

\logpage{[ 0, 0, 0 ]}
\begin{titlepage}
\mbox{}\vfill

\begin{center}{\maintitlesize \textbf{ nofoma \\
\mbox{}}}\\
\vfill

\hypersetup{pdftitle={ nofoma }}
\markright{\scriptsize \mbox{}\hfill  nofoma  \hfill\mbox{}}
{\Huge \textbf{ Normal forms of matrices \\
\mbox{}}}\\
\vfill

{\Huge  1.0.1 \mbox{}}\\[1cm]
{ 17 May 2026 \mbox{}}\\[1cm]
\mbox{}\\[2cm]
{\Large \textbf{\strut  Meinolf Geck     \strut\mbox{}}}\\
{\Large \textbf{\strut  Alia Bonnet   \strut\mbox{}}}\\
\hypersetup{pdfauthor={ Meinolf Geck    ;  Alia Bonnet  }}
\end{center}\vfill

\mbox{}\\
{\mbox{}\\
\small \noindent \textbf{ Meinolf Geck    }  Email: \href{mailto://meinolf.geck@mathematik.uni-stuttgart.de} {\texttt{meinolf.geck@mathematik.uni\texttt{\symbol{45}}stuttgart.de}}\\
  Homepage: \href{https://pnp.mathematik.uni-stuttgart.de/idsr/idsr1/geckmf/} {\texttt{https://pnp.mathematik.uni\texttt{\symbol{45}}stuttgart.de/idsr/idsr1/geckmf/}}\\
  Address: \begin{minipage}[t]{8cm}\noindent
 Fachbereich Mathematik, Pfaffenwaldring 57, 70569 Stuttgart, Germany\\
 \end{minipage}
}\\
{\mbox{}\\
\small \noindent \textbf{ Alia Bonnet  }  Email: \href{mailto://alia.bonnet@rwth.aachen.de} {\texttt{alia.bonnet@rwth.aachen.de}}}\\
\end{titlepage}

\newpage\setcounter{page}{2}
{\small 
\section*{Abstract}
\logpage{[ 0, 0, 1 ]}
 This package computes the Frobenius normal form and the
Jordan\texttt{\symbol{45}}Chevalley decomposition of a (square) matrix over
any field that is available in GAP. It also computes the Jordan normal form of
matrices over finite fields. 

 Bug reports, suggestions and comments are, welcome. Please submit them to our
issue tracker \href{https://github.com/gap-packages/nofoma/issues} {\texttt{https://github.com/gap\texttt{\symbol{45}}packages/nofoma/issues}}. \mbox{}}\\[1cm]
{\small 
\section*{Copyright}
\logpage{[ 0, 0, 2 ]}
 {\copyright} 2026 by Meinolf Geck. The \textsf{nofoma} package is free software; you can redistribute it and/or modify it under the
terms of the GNU General Public License as published by the Free Software
Foundation; either version 2 of the License, or (at your option) any later
version. \mbox{}}\\[1cm]
{\small 
\section*{Acknowledgements}
\logpage{[ 0, 0, 3 ]}
 The original code by Meinolf Geck was cleaned up and made ready for inclusion
in the \textsf{GAP} package distribution Alia Bonnet in 2025, with some assistance by Max Horn. \mbox{}}\\[1cm]
\newpage

\def\contentsname{Contents\logpage{[ 0, 0, 4 ]}}

\tableofcontents
\newpage

     
\chapter{\textcolor{Chapter }{The nofoma package}}\label{Chapter_The_nofoma_package}
\logpage{[ 1, 0, 0 ]}
\hyperdef{L}{X7D9CADDE7991BC82}{}
{
  Let $K$ be a field and $A$ be an $n\times n$\texttt{\symbol{45}}matrix over $K$. This package provides functions for computing both the Frobenius normal form
and the Jordan normal form of $A$. Furthermore, it also includes a functions for the computation of a primary
decomposition and the Jordan\texttt{\symbol{45}}Chevalley decomposition of $A$. 

 
\section{\textcolor{Chapter }{Installation of the \textsf{nofoma} package}}\label{Chapter_The_nofoma_package_Section_Installation_of_the_nofoma_package}
\logpage{[ 1, 1, 0 ]}
\hyperdef{L}{X81239C057F616FC3}{}
{
  

 To install this package first unpack it inside some \textsf{GAP} root directory in the subdirectory \texttt{pkg} (see  (\textbf{Reference: Installing a GAP Package})). Then \textsf{nofoma} can already be loaded and used (just type \texttt{LoadPackage("nofoma");}). 

 }

 }

   
\chapter{\textcolor{Chapter }{Normal forms of matrices}}\label{Chapter_Normal_forms_of_matrices}
\logpage{[ 2, 0, 0 ]}
\hyperdef{L}{X7CB17C628017D2B2}{}
{
  
\section{\textcolor{Chapter }{The Frobenius normal form}}\label{Chapter_Normal_forms_of_matrices_Section_The_Frobenius_normal_form}
\logpage{[ 2, 1, 0 ]}
\hyperdef{L}{X872D4005810328D8}{}
{
  Given a field $K$ and an $n\times n$\texttt{\symbol{45}}matrix $A$ over $K$, the \emph{Frobenius normal form} of $A$ is a block diagonal matrix, where the diagonal blocks are companion matrices
corresponding to the invariant factors of $A$. It reflects the minimal decomposition of the vector space $K^n$ into cyclic subspaces under the action of $A$. The Frobenius normal form is also called the rational canonical form. 

\subsection{\textcolor{Chapter }{FrobeniusNormalForm}}
\logpage{[ 2, 1, 1 ]}\nobreak
\hyperdef{L}{X831DC8FD82FB92CE}{}
{\noindent\textcolor{FuncColor}{$\triangleright$\enspace\texttt{FrobeniusNormalForm({\mdseries\slshape A})\index{FrobeniusNormalForm@\texttt{FrobeniusNormalForm}}
\label{FrobeniusNormalForm}
}\hfill{\scriptsize (function)}}\\


 Returns the invariant factors of a matrix \mbox{\texttt{\mdseries\slshape A}} and an invertible matrix $P$ such that $PAP^{-1}$ is the Frobenius normal form of \mbox{\texttt{\mdseries\slshape A}}. The algorithm first computes a maximal vector and an \mbox{\texttt{\mdseries\slshape A}}\texttt{\symbol{45}}invariant complement following Jacob's construction (as
described in matrix language in \cite{Gec20}); then the algorithm continues recursively. It works for matrices over any
field that is available in \textsf{GAP}. The output is a triple with 
\begin{itemize}
\item  1st component = list of invariant factors; 
\item  2nd component = base change matrix $P$; and 
\item  3rd component = indices where the various blocks in the normal form begin. 
\end{itemize}
 

 
\begin{Verbatim}[commandchars=!@|,fontsize=\small,frame=single,label=Example]
  !gapprompt@gap>| !gapinput@A:=[ [  2,  2,  0,  1,  0,  2,  1 ],|
  !gapprompt@>| !gapinput@        [  0,  4,  0,  0,  0,  1,  0 ],|
  !gapprompt@>| !gapinput@        [  0,  1,  1,  0,  0,  1,  1 ],|
  !gapprompt@>| !gapinput@        [  0, -1,  0,  1,  0, -1,  0 ],|
  !gapprompt@>| !gapinput@        [  0, -7,  0,  0,  1, -5,  0 ],|
  !gapprompt@>| !gapinput@        [  0, -2,  0,  0,  0,  1,  0 ],|
  !gapprompt@>| !gapinput@        [  0, -1,  0,  0,  0, -1,  1 ] ];;|
  !gapprompt@gap>| !gapinput@f:=FrobeniusNormalForm(A);|
  [ [ x_1^4-7*x_1^3+17*x_1^2-17*x_1+6, x_1^2-3*x_1+2, x_1-1 ], 
    [ [    1,   -2,    1,    1,    0,    0,    1 ],
      [    2,   -7,    1,    2,    0,   -1,    3 ],
      [    4,  -26,    1,    4,    0,   -8,    6 ],
      [    8,  -89,    1,    8,    0,  -35,   11 ],
      [ -1/2,   -2,    0,  1/2,    0,   -2, -3/2 ],
      [   -1,   -4,    0,    0,    0,   -4,   -2 ],
      [    0,  9/4,    0,   -3,    1,  5/4,  1/4 ] ],
    [ 1, 5, 7 ]  ]                 
  !gapprompt@gap>| !gapinput@PrintArray(f[2]*A*f[2]^-1);|
  [ [   0,   1,   0,   0,   0,   0,   0 ], 
    [   0,   0,   1,   0,   0,   0,   0 ],
    [   0,   0,   0,   1,   0,   0,   0 ],
    [  -6,  17, -17,   7,   0,   0,   0 ],
    [   0,   0,   0,   0,   0,   1,   0 ],
    [   0,   0,   0,   0,  -2,   3,   0 ],
    [   0,   0,   0,   0,   0,   0,   1 ] ]
\end{Verbatim}
 Note that the Frobenius normal form is unique up to the choice of the
companion matrices and the permutation of the blocks corresponding to the
invariant factors. So while this function is significantly more efficient than
the existing 'RationalCanonicalFormTransform', the two functions yield
slightly different results. In 'RationalCanonicalFormTransform', the companion
matrices are consistent with the output of 'CompanionMat'. However, given an $n\times n$ cyclic matrix $A$, along with a corresponding cyclic vector $v$, one can compute a change of basis matrix from A to a companion matrix of its
minimal polynomial by computing $v$ multiplied with powers of $A$ (i.e., $v$, $vA$, ...., $vA^{n-1}$). This approach follows GAP{\textquoteright}s convention of right
multiplication and yields a companion matrix in the form used by \texttt{FrobeniusNormalForm}. Furthermore 'RationalCanonicalFormTransform' sorts the invariant factors in
ascending order, while \texttt{FrobeniusNormalForm} sorts them in descending order. 
\begin{Verbatim}[commandchars=!@|,fontsize=\small,frame=single,label=Example]
  !gapprompt@gap>| !gapinput@A := [ [ 0*Z(5), Z(5)^3, 0*Z(5), Z(5)^0, Z(5)^3 ], |
  !gapprompt@>| !gapinput@  [ Z(5)^0, 0*Z(5), Z(5)^2, Z(5)^2, Z(5)^2 ], |
  !gapprompt@>| !gapinput@  [ Z(5), Z(5), 0*Z(5), Z(5)^0, Z(5)^3 ], |
  !gapprompt@>| !gapinput@  [ 0*Z(5), Z(5), Z(5)^2, Z(5)^2, Z(5) ], |
  !gapprompt@>| !gapinput@  [ Z(5)^3, Z(5)^3, Z(5)^0, Z(5)^3, Z(5)^0 ] ];;|
  !gapprompt@gap>| !gapinput@T:=RationalCanonicalFormTransform(A);;|
  !gapprompt@gap>| !gapinput@S:=TransposedMat(FrobeniusNormalForm(TransposedMat(A))[2]);;|
  !gapprompt@gap>| !gapinput@Display(A^T);|
   . . . . 3
   1 . . . 2
   . 1 . . 3
   . . 1 . 4
   . . . 1 .
  !gapprompt@gap>| !gapinput@Display(A^S);|
   . . . . 3
   1 . . . 2
   . 1 . . 3
   . . 1 . 4
   . . . 1 .
\end{Verbatim}
 

 Additionally, \texttt{RationalCanonicalFormTransform} sorts the invariant factors in ascending order, whereas the \texttt{FrobeniusNormalForm} sorts them in descending order. Consequently, the outputs of the two functions
agree up to a permutation of blocks and transposition. To get a drop in
replacement for 'RationalCanonicalFormTransform', see \texttt{FrobeniusNormalFormLikeRCFT} (\ref{FrobeniusNormalFormLikeRCFT}). 
\begin{Verbatim}[commandchars=!@|,fontsize=\small,frame=single,label=Example]
  !gapprompt@gap>| !gapinput@aa:=[[  0, -8, 12, 40,-36,  4,  0, 59, 15, -9],|
  !gapprompt@>| !gapinput@        [ -2, -2, -2,  6,-11,  1, -1, 10,  1,  0],|
  !gapprompt@>| !gapinput@        [  1,  5,  0, -6, 12, -2,  0,-12, -4,  2],|
  !gapprompt@>| !gapinput@        [  0,  0,  0,  2,  0,  0,  0,  7,  0,  0],|
  !gapprompt@>| !gapinput@        [  0,  2, -3, -7,  8, -1,  0, -7, -3,  2],|
  !gapprompt@>| !gapinput@        [ -5, -4, -6, 18,-30,  2, -2, 35,  5, -1],|
  !gapprompt@>| !gapinput@        [ -1, -6,  6, 20,-28,  3,  0, 24, 10, -6],|
  !gapprompt@>| !gapinput@        [  0,  0,  0, -1,  0,  0,  0, -3,  0,  0],|
  !gapprompt@>| !gapinput@        [  0,  0, -1, -2, -2,  0, -1, -7,  0,  0],|
  !gapprompt@>| !gapinput@        [  0, -8,  9, 21,-36,  4, -2, 12, 12, -8]];;|
  !gapprompt@gap>| !gapinput@t:=RationalCanonicalFormTransform(aa);;|
  !gapprompt@gap>| !gapinput@Display(aa^t);|
  [ [   0,   0,   0,   1,   0,   0,   0,   0,   0,   0 ],
    [   1,   0,   0,   0,   0,   0,   0,   0,   0,   0 ],
    [   0,   1,   0,   0,   0,   0,   0,   0,   0,   0 ],
    [   0,   0,   1,   0,   0,   0,   0,   0,   0,   0 ],
    [   0,   0,   0,   0,   0,   0,   0,   0,   0,   1 ],
    [   0,   0,   0,   0,   1,   0,   0,   0,   0,   1 ],
    [   0,   0,   0,   0,   0,   1,   0,   0,   0,   1 ],
    [   0,   0,   0,   0,   0,   0,   1,   0,   0,   0 ],
    [   0,   0,   0,   0,   0,   0,   0,   1,   0,  -1 ],
    [   0,   0,   0,   0,   0,   0,   0,   0,   1,  -1 ] ]
  !gapprompt@gap>| !gapinput@res:=FrobeniusNormalForm(aa);;|
  !gapprompt@gap>| !gapinput@Display(aa^(res[2]^-1));|
  [ [   0,   1,   0,   0,   0,   0,   0,   0,   0,   0 ],
    [   0,   0,   1,   0,   0,   0,   0,   0,   0,   0 ],
    [   0,   0,   0,   1,   0,   0,   0,   0,   0,   0 ],
    [   0,   0,   0,   0,   1,   0,   0,   0,   0,   0 ],
    [   0,   0,   0,   0,   0,   1,   0,   0,   0,   0 ],
    [   1,   1,   1,   0,  -1,  -1,   0,   0,   0,   0 ],
    [   0,   0,   0,   0,   0,   0,   0,   1,   0,   0 ],
    [   0,   0,   0,   0,   0,   0,   0,   0,   1,   0 ],
    [   0,   0,   0,   0,   0,   0,   0,   0,   0,   1 ],
    [   0,   0,   0,   0,   0,   0,   1,   0,   0,   0 ] ]
\end{Verbatim}
 }

 

\subsection{\textcolor{Chapter }{FrobeniusNormalFormLikeRCFT}}
\logpage{[ 2, 1, 2 ]}\nobreak
\hyperdef{L}{X81D3D6CC7E704CE7}{}
{\noindent\textcolor{FuncColor}{$\triangleright$\enspace\texttt{FrobeniusNormalFormLikeRCFT({\mdseries\slshape A})\index{FrobeniusNormalFormLikeRCFT@\texttt{FrobeniusNormalFormLikeRCFT}}
\label{FrobeniusNormalFormLikeRCFT}
}\hfill{\scriptsize (function)}}\\


 This function returns the same result as \texttt{FrobeniusNormalForm} (\ref{FrobeniusNormalForm}), except that the invariant factors are sorted in descending order and the
companion matrices on the diagonal are transposed. Furthermore, if $P$ is the computed base change matrix, the Frobenius normal form is obtained by $P^{-1}AP$ (instead of $PAP^{-1}$). This means that $P$ can be used as a direct drop\texttt{\symbol{45}}in replacement for
RationalCanonicalFormTransform. 

 Note that this function works by calling \texttt{FrobeniusNormalForm} (\ref{FrobeniusNormalForm}) and then modifying the computed transformation matrix and is thus potentially
less efficient than using the original function. 
\begin{Verbatim}[commandchars=!@|,fontsize=\small,frame=single,label=Example]
  !gapprompt@gap>| !gapinput@A := [ [ 0*Z(5), Z(5)^2, Z(5)^2, 0*Z(5), Z(5)^3, Z(5)^3, 0*Z(5), Z(5)^3, |
  !gapprompt@>| !gapinput@      0*Z(5), Z(5) ], |
  !gapprompt@>| !gapinput@  [ Z(5)^0, Z(5)^0, Z(5)^0, Z(5), Z(5)^3, Z(5), Z(5)^3, 0*Z(5), Z(5), |
  !gapprompt@>| !gapinput@      0*Z(5) ], |
  !gapprompt@>| !gapinput@  [ Z(5), 0*Z(5), 0*Z(5), Z(5)^3, Z(5)^2, Z(5)^0, 0*Z(5), Z(5)^0, |
  !gapprompt@>| !gapinput@      Z(5), Z(5)^2 ], |
  !gapprompt@>| !gapinput@  [ 0*Z(5), 0*Z(5), 0*Z(5), Z(5)^2, 0*Z(5), 0*Z(5), 0*Z(5), Z(5)^2, |
  !gapprompt@>| !gapinput@      0*Z(5), 0*Z(5) ], |
  !gapprompt@>| !gapinput@  [ 0*Z(5), Z(5)^2, Z(5)^2, Z(5)^0, Z(5)^0, Z(5)^3, 0*Z(5), Z(5)^0, |
  !gapprompt@>| !gapinput@      Z(5)^2, Z(5)^2 ], |
  !gapprompt@>| !gapinput@  [ 0*Z(5), Z(5), Z(5)^3, Z(5)^0, 0*Z(5), Z(5)^2, Z(5)^0, 0*Z(5), |
  !gapprompt@>| !gapinput@      0*Z(5), Z(5)^3 ], |
  !gapprompt@>| !gapinput@  [ Z(5)^3, Z(5)^3, Z(5), 0*Z(5), Z(5)^2, Z(5)^0, 0*Z(5), Z(5)^3, |
  !gapprompt@>| !gapinput@      0*Z(5), Z(5)^3 ], |
  !gapprompt@>| !gapinput@  [ 0*Z(5), 0*Z(5), 0*Z(5), Z(5)^3, 0*Z(5), 0*Z(5), 0*Z(5), Z(5)^2, |
  !gapprompt@>| !gapinput@      0*Z(5), 0*Z(5) ], |
  !gapprompt@>| !gapinput@  [ 0*Z(5), 0*Z(5), Z(5)^3, Z(5)^0, Z(5)^0, 0*Z(5), Z(5)^3, Z(5)^0, |
  !gapprompt@>| !gapinput@      0*Z(5), 0*Z(5) ], |
  !gapprompt@>| !gapinput@  [ 0*Z(5), Z(5)^2, Z(5)^3, Z(5), Z(5)^3, Z(5)^3, Z(5)^0, Z(5)^2, |
  !gapprompt@>| !gapinput@      Z(5)^2, Z(5)^2 ] ];;|
  !gapprompt@gap>| !gapinput@frob := FrobeniusNormalFormLikeRCFT(A)[2];;|
  !gapprompt@gap>| !gapinput@rat := RationalCanonicalFormTransform(A);;|
  !gapprompt@gap>| !gapinput@Display(A^rat);|
   . . . 1 . . . . . .
   1 . . . . . . . . .
   . 1 . . . . . . . .
   . . 1 . . . . . . .
   . . . . . . . . . 4
   . . . . 1 . . . . 2
   . . . . . 1 . . . 1
   . . . . . . 1 . . .
   . . . . . . . 1 . 1
   . . . . . . . . 1 3
  !gapprompt@gap>| !gapinput@Display(A^frob);|
   . . . 1 . . . . . .
   1 . . . . . . . . .
   . 1 . . . . . . . .
   . . 1 . . . . . . .
   . . . . . . . . . 4
   . . . . 1 . . . . 2
   . . . . . 1 . . . 1
   . . . . . . 1 . . .
   . . . . . . . 1 . 1
   . . . . . . . . 1 3
\end{Verbatim}
 }

 

\subsection{\textcolor{Chapter }{InvariantFactorsMat}}
\logpage{[ 2, 1, 3 ]}\nobreak
\hyperdef{L}{X7B675811827CBCEC}{}
{\noindent\textcolor{FuncColor}{$\triangleright$\enspace\texttt{InvariantFactorsMat({\mdseries\slshape A})\index{InvariantFactorsMat@\texttt{InvariantFactorsMat}}
\label{InvariantFactorsMat}
}\hfill{\scriptsize (function)}}\\


 Returns the invariant factors of the matrix \mbox{\texttt{\mdseries\slshape A}}, i.e., the minimal polynomials of the diagonal blocks in the Frobenius normal
form of \mbox{\texttt{\mdseries\slshape A}}. Thus, 'InvariantFactorsMat' also specifies the rational canonical form of \mbox{\texttt{\mdseries\slshape A}}, but without computing the base change. 

 
\begin{Verbatim}[commandchars=!@|,fontsize=\small,frame=single,label=Example]
  !gapprompt@gap>| !gapinput@A := [ [ 2,  2, 0, 1, 0,  2, 1 ],|
  !gapprompt@>| !gapinput@          [ 0,  4, 0, 0, 0,  1, 0 ],|
  !gapprompt@>| !gapinput@          [ 0,  1, 1, 0, 0,  1, 1 ],|
  !gapprompt@>| !gapinput@          [ 0, -1, 0, 1, 0, -1, 0 ],|
  !gapprompt@>| !gapinput@          [ 0, -7, 0, 0, 1, -5, 0 ],|
  !gapprompt@>| !gapinput@          [ 0, -2, 0, 0, 0,  1, 0 ],|
  !gapprompt@>| !gapinput@          [ 0, -1, 0, 0, 0, -1, 1 ] ];;|
  !gapprompt@gap>| !gapinput@InvariantFactorsMat(A);|
    [ x_1^4-7*x_1^3+17*x_1^2-17*x_1+6, x_1^2-3*x_1+2, x_1-1 ]
\end{Verbatim}
 }

 }

 
\section{\textcolor{Chapter }{The Jordan normal form}}\label{Chapter_Normal_forms_of_matrices_Section_The_Jordan_normal_form}
\logpage{[ 2, 2, 0 ]}
\hyperdef{L}{X8617E70787358AAD}{}
{
  The Jordan normal form of a matrix $A$ is a block diagonal matrix, where the diagonal blocks are Jordan blocks
corresponding to the elementary divisors of $A$. It reflects the maximal decomposition of the vector space $K^n$ into cyclic subspaces under the action of $A$. For a more thorough definition of the Jordan normal form and details about
the algorithms used, see \cite{Bon26}. 

\subsection{\textcolor{Chapter }{JordanNormalForm}}
\logpage{[ 2, 2, 1 ]}\nobreak
\hyperdef{L}{X86D2B8A98222AC4C}{}
{\noindent\textcolor{FuncColor}{$\triangleright$\enspace\texttt{JordanNormalForm({\mdseries\slshape A})\index{JordanNormalForm@\texttt{JordanNormalForm}}
\label{JordanNormalForm}
}\hfill{\scriptsize (function)}}\\


 Returns a list containing two entries. The first is a base change matrix $B$ such that $B$\mbox{\texttt{\mdseries\slshape A}}$B^{-1}$ is the Jordan normal form of \mbox{\texttt{\mdseries\slshape A}}, i.e. a block diagonal matrix where the diagonal blocks are Jordan blocks
corresponding to the elementary divisors of \mbox{\texttt{\mdseries\slshape A}} in descending order. The second entry is a list of the elementary divisors of \mbox{\texttt{\mdseries\slshape A}}, also in descending order. The algorithm first computes a primary
decomposition of \mbox{\texttt{\mdseries\slshape A}} and then computes a cyclic decomposition of the primary components. Finally it
computes Jordan block form for each of the cyclic components. The blocks are
ordered in the same order as in the list containing the elementary divisors.
It works for matrices over finite fields. 

 Since all of the blocks on the resulting transformed matrix are cyclic, one
can retrieve their size by the degrees of the respective elementary divisor. 

 
\begin{Verbatim}[commandchars=!@|,fontsize=\small,frame=single,label=Example]
  !gapprompt@gap>| !gapinput@A := [ [ 0*Z(5), 0*Z(5), Z(5)^3, Z(5)^3, Z(5)^3, Z(5)^0 ], |
  !gapprompt@>| !gapinput@   [ 0*Z(5), Z(5)^2, Z(5)^2, Z(5)^0, Z(5)^3, Z(5)^3 ], |
  !gapprompt@>| !gapinput@   [ Z(5)^0, Z(5)^0, Z(5)^3, Z(5)^2, Z(5)^0, Z(5) ], |
  !gapprompt@>| !gapinput@   [ 0*Z(5), Z(5)^3, Z(5), Z(5), 0*Z(5), Z(5)^2 ], |
  !gapprompt@>| !gapinput@   [ Z(5)^2, Z(5)^0, Z(5)^0, 0*Z(5), Z(5), Z(5) ], |
  !gapprompt@>| !gapinput@   [ 0*Z(5), Z(5)^0, Z(5)^2, Z(5), Z(5), Z(5) ] ];;|
  !gapprompt@gap>| !gapinput@B := JordanNormalForm(A);;|
  !gapprompt@gap>| !gapinput@Display(A^Inverse(B[1]));|
  3 . . . . .
  . 1 . . . .
  . . . 1 . .
  . . 2 . . .
  . . . . . 1
  . . . . 3 4
\end{Verbatim}
 This function computes the Jordan normal form of $A$ significantly faster if $A$ is either cyclic or has irreducible minimal polynomial. 
\begin{Verbatim}[commandchars=!@|,fontsize=\small,frame=single,label=Example]
  !gapprompt@gap>| !gapinput@B:= [ [ Z(5), Z(5)^3, 0*Z(5), 0*Z(5), 0*Z(5), 0*Z(5), 0*Z(5), 0*Z(5), 0*Z(5), 0*Z(5) ], |
  !gapprompt@>| !gapinput@[ Z(5)^0, Z(5), 0*Z(5), 0*Z(5), 0*Z(5), 0*Z(5), 0*Z(5), 0*Z(5), 0*Z(5), 0*Z(5) ], |
  !gapprompt@>| !gapinput@[ 0*Z(5), 0*Z(5), Z(5), Z(5)^3, 0*Z(5), 0*Z(5), 0*Z(5), 0*Z(5), 0*Z(5), 0*Z(5) ], |
  !gapprompt@>| !gapinput@[ 0*Z(5), 0*Z(5), Z(5)^0, Z(5), 0*Z(5), 0*Z(5), 0*Z(5), 0*Z(5), 0*Z(5), 0*Z(5) ], |
  !gapprompt@>| !gapinput@[ 0*Z(5), 0*Z(5), 0*Z(5), 0*Z(5), Z(5), Z(5)^3, 0*Z(5), 0*Z(5), 0*Z(5), 0*Z(5) ], |
  !gapprompt@>| !gapinput@[ 0*Z(5), 0*Z(5), 0*Z(5), 0*Z(5), Z(5)^0, Z(5), 0*Z(5), 0*Z(5), 0*Z(5), 0*Z(5) ], |
  !gapprompt@>| !gapinput@[ 0*Z(5), 0*Z(5), 0*Z(5), 0*Z(5), 0*Z(5), 0*Z(5), Z(5), Z(5)^3, 0*Z(5), 0*Z(5) ], |
  !gapprompt@>| !gapinput@[ 0*Z(5), 0*Z(5), 0*Z(5), 0*Z(5), 0*Z(5), 0*Z(5), Z(5)^0, Z(5), 0*Z(5), 0*Z(5) ], |
  !gapprompt@>| !gapinput@[ 0*Z(5), 0*Z(5), 0*Z(5), 0*Z(5), 0*Z(5), 0*Z(5), 0*Z(5), 0*Z(5), Z(5), Z(5)^3 ], |
  !gapprompt@>| !gapinput@[ 0*Z(5), 0*Z(5), 0*Z(5), 0*Z(5), 0*Z(5), 0*Z(5), 0*Z(5), 0*Z(5), Z(5)^0, Z(5) ] ];;|
  !gapprompt@gap>| !gapinput@Factors(MinimalPolynomial(B));|
  [ x_1^2+x_1+Z(5)^0 ]
  !gapprompt@gap>| !gapinput@Display(B^Inverse(JordanNormalForm(B)[1]));|
  . 1 . . . . . . . .
  4 4 . . . . . . . .
  . . . 1 . . . . . .
  . . 4 4 . . . . . .
  . . . . . 1 . . . .
  . . . . 4 4 . . . .
  . . . . . . . 1 . .
  . . . . . . 4 4 . .
  . . . . . . . . . 1
  . . . . . . . . 4 4
\end{Verbatim}
 }

 }

 }

   
\chapter{\textcolor{Chapter }{Other functionality}}\label{Chapter_Other_functionality}
\logpage{[ 3, 0, 0 ]}
\hyperdef{L}{X86888B667BC469D8}{}
{
  
\section{\textcolor{Chapter }{The Jordan\texttt{\symbol{45}}Chevalley decomposition}}\label{Chapter_Other_functionality_Section_The_Jordan-Chevalley_decomposition}
\logpage{[ 3, 1, 0 ]}
\hyperdef{L}{X85FD78337E4F199D}{}
{
  

\subsection{\textcolor{Chapter }{JordanChevalleyDecMat}}
\logpage{[ 3, 1, 1 ]}\nobreak
\hyperdef{L}{X7E20260B7DE0AAA2}{}
{\noindent\textcolor{FuncColor}{$\triangleright$\enspace\texttt{JordanChevalleyDecMat({\mdseries\slshape A, f})\index{JordanChevalleyDecMat@\texttt{JordanChevalleyDecMat}}
\label{JordanChevalleyDecMat}
}\hfill{\scriptsize (function)}}\\


 Returns the unique pair of matrices $D$, $N$ such that the matrix \mbox{\texttt{\mdseries\slshape A}} is written as $A=D+N$, where $N$ is a nilpotent matrix and $D$ is a matrix that is diagonalisable (over some extension field of the default
field of \mbox{\texttt{\mdseries\slshape A}}), such that $D.N=N.D$; the argument \mbox{\texttt{\mdseries\slshape f}} is a polynomial such that $f(A)=0$ (e.g., the minimal polynomial of \mbox{\texttt{\mdseries\slshape A}}). This is called the Jordan\texttt{\symbol{45}}Chevalley decomposition of \mbox{\texttt{\mdseries\slshape A}}; the algorithm is based on \cite{Gec22}. Note that this algorithm does not require the knowledge of the eigenvalues
of \mbox{\texttt{\mdseries\slshape A}}; it works over any perfect field that is available in \textsf{GAP}. 

 
\begin{Verbatim}[commandchars=!@|,fontsize=\small,frame=single,label=Example]
  !gapprompt@gap>| !gapinput@A:=[ [  6, -2,  6,  1,  1 ],|
  !gapprompt@>| !gapinput@        [  1, -1,  2,  1, -2 ],|
  !gapprompt@>| !gapinput@        [ -2,  0, -1,  0, -1 ],|
  !gapprompt@>| !gapinput@        [ -1,  0, -2,  2, -1 ],|
  !gapprompt@>| !gapinput@        [ -4,  4, -6, -2,  3 ] ];;|
  !gapprompt@gap>| !gapinput@jc:=JordanChevalleyDecMat(A,MinimalPolynomial(A));|
  [ [ [  4,  0,  4, -1,  1 ], 
      [  1,  0,  1,  1, -1 ], 
      [ -1, -1,  0,  1, -1 ], 
      [  0,  0, -2,  3,  0 ], 
      [ -3,  2, -4, -1,  2 ] ], 
    [ [  2, -2,  2,  2,  0 ], 
      [  0, -1,  1,  0, -1 ], 
      [ -1,  1, -1, -1,  0 ], 
      [ -1,  0,  0, -1, -1 ], 
      [ -1,  2, -2, -1,  1 ] ] ]
  !gapprompt@gap>| !gapinput@MinimalPolynomial(jc[1]);|
  x_1^3-5*x_1^2+9*x_1-5
  !gapprompt@gap>| !gapinput@Factors(last);|
  [ x_1-1, x_1^2-4*x_1+5 ]  
  !gapprompt@gap>| !gapinput@MinimalPolynomial(jc[2]);|
  x_1^2                     
\end{Verbatim}
 If the input matrix is very large, then \texttt{JordanChevalleyDecMatF} (\ref{JordanChevalleyDecMatF}) may be more efficient; this function first computes the Frobenius normal form
of \mbox{\texttt{\mdseries\slshape A}} and then applies \texttt{JordanChevalleyDecMat} to each diagonal block. (The result will be the same as that of
'JordanChevalleyDecMat(\mbox{\texttt{\mdseries\slshape A}});)' }

 

\subsection{\textcolor{Chapter }{JordanChevalleyDecMatF}}
\logpage{[ 3, 1, 2 ]}\nobreak
\hyperdef{L}{X811E309378DAD798}{}
{\noindent\textcolor{FuncColor}{$\triangleright$\enspace\texttt{JordanChevalleyDecMatF({\mdseries\slshape A})\index{JordanChevalleyDecMatF@\texttt{JordanChevalleyDecMatF}}
\label{JordanChevalleyDecMatF}
}\hfill{\scriptsize (function)}}\\


 First computes the Frobenius normal form and then applies \texttt{JordanChevalleyDecMat} (\ref{JordanChevalleyDecMat}) to each diagonal block. }

 }

 
\section{\textcolor{Chapter }{The primary decomposition}}\label{Chapter_Other_functionality_Section_The_primary_decomposition}
\logpage{[ 3, 2, 0 ]}
\hyperdef{L}{X8202AE3378FA1965}{}
{
  

\subsection{\textcolor{Chapter }{PrimaryDecomposition (for IsMatrixOrMatrixObj)}}
\logpage{[ 3, 2, 1 ]}\nobreak
\hyperdef{L}{X7DD44F8E79BFDE8B}{}
{\noindent\textcolor{FuncColor}{$\triangleright$\enspace\texttt{PrimaryDecomposition({\mdseries\slshape A})\index{PrimaryDecomposition@\texttt{PrimaryDecomposition}!for IsMatrixOrMatrixObj}
\label{PrimaryDecomposition:for IsMatrixOrMatrixObj}
}\hfill{\scriptsize (attribute)}}\\


 Returns a list containing three elements. The first element is a base change
matrix $B$ such that $B$\mbox{\texttt{\mdseries\slshape A}}$B^{-1}$ is a primary form of the matrix \mbox{\texttt{\mdseries\slshape A}}, i.e., a block diagonal matrix where the minimal polynomials of the the
diagonal blocks are precisely the powers of irreducible factors of the minimal
polynomial of \mbox{\texttt{\mdseries\slshape A}}, in descending order. The second element is a list containing the collected
irreducible factors of the minimal polynomial of \mbox{\texttt{\mdseries\slshape A}}, in the same order. The last element is a list containing the the size of
each block. The exact algorithm used in this function is described in \cite{Bon26} 

 
\begin{Verbatim}[commandchars=!@|,fontsize=\small,frame=single,label=Example]
  !gapprompt@gap>| !gapinput@A := [ [ Z(5)^2, 0*Z(5), Z(5)^2, Z(5)^3, Z(5) ], |
  !gapprompt@>| !gapinput@   [ 0*Z(5), 0*Z(5), Z(5)^3, Z(5), Z(5)^0 ],  |
  !gapprompt@>| !gapinput@   [ Z(5), Z(5)^0, 0*Z(5), Z(5)^0, 0*Z(5) ],|
  !gapprompt@>| !gapinput@   [ Z(5)^0, Z(5)^0, Z(5)^0, 0*Z(5), Z(5)^3 ],|
  !gapprompt@>| !gapinput@   [ Z(5), 0*Z(5), Z(5)^3, 0*Z(5), Z(5)^3 ] ];;|
  !gapprompt@gap>| !gapinput@B := PrimaryDecomposition(A);;|
  !gapprompt@gap>| !gapinput@Display(B[3]);|
  [ 1, 4 ]
  !gapprompt@gap>| !gapinput@Factors(MinimalPolynomial(A));|
  [ x_1-Z(5)^0, x_1^4-x_1^3+Z(5)^3*x_1+Z(5)^3 ]
  !gapprompt@gap>| !gapinput@PrimA := A^Inverse(B[1]);;|
  !gapprompt@gap>| !gapinput@MinimalPolynomial(PrimA{[1..1]}{[1..1]});|
  x_1-Z(5)^0
  !gapprompt@gap>| !gapinput@MinimalPolynomial(PrimA{[2..5]}{[2..5]});|
  x_1^4-x_1^3+Z(5)^3*x_1+Z(5)^3
\end{Verbatim}
 }

 }

 }

   
\chapter{\textcolor{Chapter }{Auxiliary functions}}\label{Chapter_Auxiliary_functions}
\logpage{[ 4, 0, 0 ]}
\hyperdef{L}{X866E18057EF83F65}{}
{
  
\section{\textcolor{Chapter }{Vectors and matrices and their associated polynomials}}\label{Chapter_Auxiliary_functions_Section_Vectors_and_matrices_and_their_associated_polynomials}
\logpage{[ 4, 1, 0 ]}
\hyperdef{L}{X78EBB30F827DBDF4}{}
{
  

\subsection{\textcolor{Chapter }{GcdCoprimeSplit}}
\logpage{[ 4, 1, 1 ]}\nobreak
\hyperdef{L}{X868E1E5878D2DB25}{}
{\noindent\textcolor{FuncColor}{$\triangleright$\enspace\texttt{GcdCoprimeSplit({\mdseries\slshape a, b})\index{GcdCoprimeSplit@\texttt{GcdCoprimeSplit}}
\label{GcdCoprimeSplit}
}\hfill{\scriptsize (function)}}\\


 Computes a divisor $a_1$ of the polynomial \mbox{\texttt{\mdseries\slshape a}} and a divisor $b_1$ of the polynomial \mbox{\texttt{\mdseries\slshape b}} such that $a_1$ and $b_1$ are coprime and the lcm of \mbox{\texttt{\mdseries\slshape a}}, \mbox{\texttt{\mdseries\slshape b}} is $a_1*b_1$. This is based on Lemma 5 in \cite{Bon14}. (See also Lemma 4.3 in \cite{Gec20}). 

 (Note that it does not use the prime factorisation of polynomials but only gcd
computations.) 

 
\begin{Verbatim}[commandchars=!@|,fontsize=\small,frame=single,label=Example]
  !gapprompt@gap>| !gapinput@x:=X(Rationals);;|
  !gapprompt@gap>| !gapinput@a:=x^2*(x-1)^3*(x-2)*(x-3);|
  x_1^7-8*x_1^6+24*x_1^5-34*x_1^4+23*x_1^3-6*x_1^2
  !gapprompt@gap>| !gapinput@b:=x^2*(x-1)^2*(x-2)^4*(x-4);|
  x_1^9-14*x_1^8+81*x_1^7-252*x_1^6+456*x_1^5-480*x_1^4+272*x_1^3-64*x_1^2
  !gapprompt@gap>| !gapinput@GcdCoprimeSplit(a,b);|
  [ x_1^5-4*x_1^4+5*x_1^3-2*x_1^2, x_1^4-6*x_1^3+12*x_1^2-10*x_1+3, 
    x_1^7-12*x_1^6+56*x_1^5-128*x_1^4+144*x_1^3-64*x_1^2 ]
\end{Verbatim}
 }

 

\subsection{\textcolor{Chapter }{PolynomialToMatVec}}
\logpage{[ 4, 1, 2 ]}\nobreak
\hyperdef{L}{X8228A5537A30AD38}{}
{\noindent\textcolor{FuncColor}{$\triangleright$\enspace\texttt{PolynomialToMatVec({\mdseries\slshape A, pol, v})\index{PolynomialToMatVec@\texttt{PolynomialToMatVec}}
\label{PolynomialToMatVec}
}\hfill{\scriptsize (function)}}\\


 Returns the row vector obtained by multiplying the row vector \mbox{\texttt{\mdseries\slshape v}} with the matrix \mbox{\texttt{\mdseries\slshape pol}}(\mbox{\texttt{\mdseries\slshape A}}), where p is a polynomial. The actual computation is more efficient than this
naive approach. 

 
\begin{Verbatim}[commandchars=!@|,fontsize=\small,frame=single,label=Example]
  !gapprompt@gap>| !gapinput@A:=[ [ 0, 1, 0, 1 ],|
  !gapprompt@>| !gapinput@        [ 0, 0, 0, 0 ],|
  !gapprompt@>| !gapinput@        [ 0, 1, 0, 1 ],|
  !gapprompt@>| !gapinput@        [ 1, 1, 1, 1 ] ];;|
  !gapprompt@gap>| !gapinput@x:=X(Rationals);;|
  !gapprompt@gap>| !gapinput@pol:=x^6-6*x^5+12*x^4-10*x^3+3*x^2;;|
  !gapprompt@gap>| !gapinput@v:=[ 1, 1, 1, 1 ];;|
  !gapprompt@gap>| !gapinput@PolynomialToMatVec(A,pol,v);|
  [ 8, -16, 8, -16 ]
\end{Verbatim}
 }

 

\subsection{\textcolor{Chapter }{LcmMaximalVectorMat}}
\logpage{[ 4, 1, 3 ]}\nobreak
\hyperdef{L}{X7D07B0E67857FFA4}{}
{\noindent\textcolor{FuncColor}{$\triangleright$\enspace\texttt{LcmMaximalVectorMat({\mdseries\slshape A, v1, v2, pol1, pol2})\index{LcmMaximalVectorMat@\texttt{LcmMaximalVectorMat}}
\label{LcmMaximalVectorMat}
}\hfill{\scriptsize (function)}}\\


 Returns, given a matrix \mbox{\texttt{\mdseries\slshape A}}, vectors \mbox{\texttt{\mdseries\slshape v1}}, \mbox{\texttt{\mdseries\slshape v2}} with minimal polynomials \mbox{\texttt{\mdseries\slshape pol1}}, \mbox{\texttt{\mdseries\slshape pol2}}, a new pair [$v$,$pol$], where $v$ has minimal polynomial $pol$, and $pol$ is the least common multiple of \mbox{\texttt{\mdseries\slshape pol1}} and \mbox{\texttt{\mdseries\slshape pol2}}. This crucially relies on \texttt{GcdCoprimeSplit} (\ref{GcdCoprimeSplit}) to avoid factorisation of polynomials. }

 

\subsection{\textcolor{Chapter }{SpinMatVector}}
\logpage{[ 4, 1, 4 ]}\nobreak
\hyperdef{L}{X84C7CEF97C03673D}{}
{\noindent\textcolor{FuncColor}{$\triangleright$\enspace\texttt{SpinMatVector({\mdseries\slshape A, v})\index{SpinMatVector@\texttt{SpinMatVector}}
\label{SpinMatVector}
}\hfill{\scriptsize (function)}}\\


 Computes the smallest subspace containing the vector \mbox{\texttt{\mdseries\slshape v}} that is invariant under the matrix \mbox{\texttt{\mdseries\slshape A}}. The output is a quadruple, with 
\begin{itemize}
\item  1st component = basis of that subspace in row echelon form; 
\item  2nd component = matrix with rows $\mbox{\texttt{\mdseries\slshape v}}, \mbox{\texttt{\mdseries\slshape v}}.\mbox{\texttt{\mdseries\slshape A}}, \mbox{\texttt{\mdseries\slshape v}}.\mbox{\texttt{\mdseries\slshape A}}^2, ..., \mbox{\texttt{\mdseries\slshape v}}.\mbox{\texttt{\mdseries\slshape A}}^{{d-1}}$ (where $d$ is the degree of the local minimal polynomial of \mbox{\texttt{\mdseries\slshape v}}); 
\item  3rd component = the coefficients $a_0$, $a_1$, ..., $a_d$ of the local minimal polynomial; and 
\item  4th component = the positions of the pivots of the first component. 
\end{itemize}
 

 
\begin{Verbatim}[commandchars=!@|,fontsize=\small,frame=single,label=Example]
  !gapprompt@gap>| !gapinput@A:=[ [   5,   2,  -4,   2 ],|
  !gapprompt@>| !gapinput@        [  -1,   0,   2,  -1 ],|
  !gapprompt@>| !gapinput@        [  -1,  -1,   3,  -1 ],|
  !gapprompt@>| !gapinput@        [ -13,  -7,  14,  -6 ] ];;|
  !gapprompt@gap>| !gapinput@SpinMatVector(A,[1,0,0,0]);|
  [ [ [ 1, 0, 0, 0 ], [ 0, 1, -2, 1 ] ],
    [ [ 1, 0, 0, 0 ], [ 5, 2, -4, 2 ] ],
    [ -1, 0, 1 ],
    [ 1, 2 ] ]
  !gapprompt@gap>| !gapinput@SpinMatVector(A,[0,1,0,0]);|
  [ [ [ 0, 1, 0, 0 ], [ 1, 0, -2, 1 ], [ 0, 0, 1, -1/2 ] ],
    [ [ 0, 1, 0, 0 ], [ -1, 0, 2, -1 ], [ 6, 3, -4, 2 ] ],
    [ 1, -1, -1, 1 ],
    [ 2, 1, 3 ] ]
  !gapprompt@gap>| !gapinput@SpinMatVector(A,[1,1,0,0]);|
  [ [ [ 1, 1, 0, 0 ], [ 0, 1, 1, -1/2 ] ],
    [ [ 1, 1, 0, 0 ], [ 4, 2, -2, 1 ] ],
    [ 1, -2, 1 ],
    [ 1, 2 ] ]
\end{Verbatim}
 }

 

\subsection{\textcolor{Chapter }{CyclicChainMat}}
\logpage{[ 4, 1, 5 ]}\nobreak
\hyperdef{L}{X810E09A77D1CA2A8}{}
{\noindent\textcolor{FuncColor}{$\triangleright$\enspace\texttt{CyclicChainMat({\mdseries\slshape A})\index{CyclicChainMat@\texttt{CyclicChainMat}}
\label{CyclicChainMat}
}\hfill{\scriptsize (function)}}\\


 Repeatedly computes the smallest invariant subspaces containing different
vectors to compute a chain of cyclic subspaces. The output is a triple \texttt{[B,C,svec]} where $C$ is such that $C\mbox{\texttt{\mdseries\slshape A}}C^{-1}$ has a block triangular shape with companion matrices along the diagonal), $B$ is the row echelon form of C and svec is the list of indices where the blocks
begin. 

 
\begin{Verbatim}[commandchars=!@|,fontsize=\small,frame=single,label=Example]
  !gapprompt@gap>| !gapinput@A:=[ [ 0, 1, 0, 1 ],|
  !gapprompt@>| !gapinput@        [ 0, 0, 1, 0 ],|
  !gapprompt@>| !gapinput@        [ 0, 1, 0, 1 ],|
  !gapprompt@>| !gapinput@        [ 1, 1, 1, 1 ] ];;|
  !gapprompt@gap>| !gapinput@sp:=CyclicChainMat(A);|
  [ [ [ 1, 0, 0, 0 ], [ 0, 1, 0, 1 ], [ 0, 0, 1, 0 ], [ 0, 0, 0, 1 ] ],
    [ [ 1, 0, 0, 0 ], [ 0, 1, 0, 1 ], [ 1, 1, 2, 1 ], [ 0, 0, 0, 1 ] ],
    [ 1, 4, 5 ] ]
  !gapprompt@gap>| !gapinput@PrintArray(sp[2]*A*sp[2]^-1);|
  [ [    0,    1,    0,    0 ],
    [    0,    0,    1,    0 ],
    [    0,    3,    1,    0 ],
    [  1/2,  1/2,  1/2,    0 ] ]
\end{Verbatim}
 }

 

\subsection{\textcolor{Chapter }{MaximalVectorMat}}
\logpage{[ 4, 1, 6 ]}\nobreak
\hyperdef{L}{X7915F1A4798392DF}{}
{\noindent\textcolor{FuncColor}{$\triangleright$\enspace\texttt{MaximalVectorMat({\mdseries\slshape A})\index{MaximalVectorMat@\texttt{MaximalVectorMat}}
\label{MaximalVectorMat}
}\hfill{\scriptsize (function)}}\\


 Returns the minimal polynomial and a maximal vector of the matrix \mbox{\texttt{\mdseries\slshape A}}, that is, a vector whose local minimal polynomial is that of \mbox{\texttt{\mdseries\slshape A}}. This is done by repeatedly spinning up vectors until a maximal one is found.
The exact algorithm is a combination of 
\begin{itemize}
\item  the minimal polynomial algorithm by Neunhoeffer\texttt{\symbol{45}}Praeger;
see \cite{Neu08}; and 
\item  the method described by Bongartz (see \cite{Bon14}) for computing maximal vectors. 
\end{itemize}
 

 See also the article by Geck at \cite{Gec20}. 

 
\begin{Verbatim}[commandchars=!@|,fontsize=\small,frame=single,label=Example]
  !gapprompt@gap>| !gapinput@A:=[ [  2,  2,  0,  1,  0,  2,  1 ],|
  !gapprompt@>| !gapinput@        [  0,  4,  0,  0,  0,  1,  0 ],|
  !gapprompt@>| !gapinput@        [  0,  1,  1,  0,  0,  1,  1 ],|
  !gapprompt@>| !gapinput@        [  0, -1,  0,  1,  0, -1,  0 ],|
  !gapprompt@>| !gapinput@        [  0, -7,  0,  0,  1, -5,  0 ],|
  !gapprompt@>| !gapinput@        [  0, -2,  0,  0,  0,  1,  0 ],|
  !gapprompt@>| !gapinput@        [  0, -1,  0,  0,  0, -1,  1 ] ];;|
  !gapprompt@gap>| !gapinput@MaximalVectorMat(A);|
  [ [ 1, -2, 1, 1, 0, 0, 1 ], x_1^4-7*x_1^3+17*x_1^2-17*x_1+6 ]
  !gapprompt@gap>| !gapinput@v:=last[1];;|
  !gapprompt@gap>| !gapinput@SpinMatVector(A,v)[3];|
  [ 6, -17, 17, -7, 1 ]
\end{Verbatim}
 In the following example, $M_2$ is the (challenging) test matrix from the paper by
Neunhoeffer\texttt{\symbol{45}}Praeger: 

 
\begin{Verbatim}[commandchars=!@|,fontsize=\small,frame=single,label=Example]
  !gapprompt@gap>| !gapinput@LoadPackage("AtlasRep");; g:=AtlasGroup("B",1); M2:=g.1+g.2+g.1*g.2;|
  <matrix group of size 4154781481226426191177580544000000 with 2 generators>
  <an immutable 4370x4370 matrix over GF2>
  !gapprompt@gap>| !gapinput@SetInfoLevel(Infonofoma,1);|
  !gapprompt@gap>| !gapinput@MaximalVectorMat(M2);;time;|
  #I Degree of minimal polynomial is 2097
  6725
  !gapprompt@gap>| !gapinput@MinimalPolynomial(M2);;time;|
  13415
  !gapprompt@gap>| !gapinput@LoadPackage("cvec");|
  !gapprompt@gap>| !gapinput@MinimalPolynomial(CMat(M2));;time;|
  9721
\end{Verbatim}
 }

 

\subsection{\textcolor{Chapter }{JacobMatComplement}}
\logpage{[ 4, 1, 7 ]}\nobreak
\hyperdef{L}{X83FABFE183BD7477}{}
{\noindent\textcolor{FuncColor}{$\triangleright$\enspace\texttt{JacobMatComplement({\mdseries\slshape T, d})\index{JacobMatComplement@\texttt{JacobMatComplement}}
\label{JacobMatComplement}
}\hfill{\scriptsize (function)}}\\


 Modifies an already given complementary subspace to the complementary subspace
defined by Jacob; concretely, this is realized by assuming that \mbox{\texttt{\mdseries\slshape T}} is a matrix in block triangular shape, where the upper left diagonal block is
a companion matrix (as returned by \texttt{RatFormStep1} (\ref{RatFormStep1}); the variable \mbox{\texttt{\mdseries\slshape d}} gives the size of that block. (If \mbox{\texttt{\mdseries\slshape T}} gives a maximal cyclic subspace, then Jacob's complement is also \mbox{\texttt{\mdseries\slshape T}}\texttt{\symbol{45}}invariant; but even if not, it appears to be very useful
because it produces many zeroes.) }

 

\subsection{\textcolor{Chapter }{RatFormStep1}}
\logpage{[ 4, 1, 8 ]}\nobreak
\hyperdef{L}{X8741791C7A022080}{}
{\noindent\textcolor{FuncColor}{$\triangleright$\enspace\texttt{RatFormStep1({\mdseries\slshape A, v})\index{RatFormStep1@\texttt{RatFormStep1}}
\label{RatFormStep1}
}\hfill{\scriptsize (function)}}\\


 Spins up a vector \mbox{\texttt{\mdseries\slshape v}} under a matrix \mbox{\texttt{\mdseries\slshape A}}, computes a complementary subspace (using Jacob's construction), and performs
the base change. The output is a quadruple \texttt{[A1,P,pol,str]} where $A1$ is the new matrix, $P$ is the base change, $pol$ is the minimal polynomial and $str$ is either 'split' or 'not', according to whether the extension is split or
not. 

 
\begin{Verbatim}[commandchars=!@|,fontsize=\small,frame=single,label=Example]
  !gapprompt@gap>| !gapinput@v:=[ 1, 1, 1, 1 ];;|
  !gapprompt@gap>| !gapinput@A:=[ [ 0, 1, 0, 1 ],|
  !gapprompt@>| !gapinput@        [ 0, 0, 1, 0 ],|
  !gapprompt@>| !gapinput@        [ 0, 1, 0, 1 ],|
  !gapprompt@>| !gapinput@        [ 1, 1, 1, 1 ] ];;|
  !gapprompt@gap>| !gapinput@PrintArray(RatFormStep1(A,v)[1]);|
  [ [  0,  1,  0,  0 ],
    [  0,  0,  1,  0 ],
    [  0,  3,  1,  0 ],
    [  1,  0,  0,  0 ] ]
\end{Verbatim}
 }

 }

 }

 \def\bibname{References\logpage{[ "Bib", 0, 0 ]}
\hyperdef{L}{X7A6F98FD85F02BFE}{}
}

\bibliographystyle{alpha}
\bibliography{nofoma.bib}

\addcontentsline{toc}{chapter}{References}

\def\indexname{Index\logpage{[ "Ind", 0, 0 ]}
\hyperdef{L}{X83A0356F839C696F}{}
}

\cleardoublepage
\phantomsection
\addcontentsline{toc}{chapter}{Index}


\printindex

\immediate\write\pagenrlog{["Ind", 0, 0], \arabic{page},}
\newpage
\immediate\write\pagenrlog{["End"], \arabic{page}];}
\immediate\closeout\pagenrlog
\end{document}
